\documentclass[10pt,UTF8]{article}

\usepackage{geometry}

\usepackage{ctex}

\geometry{a4paper,scale=0.9}
\ctexset{today=old}

\usepackage[bookmarksnumbered=true]{hyperref}  % 生成大纲


\usepackage{times}  % 设置正文字体为 Adobe Times Roman
% \usepackage{mathptmx} % 只加载数学字体
% \renewcommand{\rmdefault}{cmr} % 恢复正文字体为 Computer Modern Roman

% \usepackage{makecell}
% \usepackage{multirow} % 表格多行合并

% \usepackage{graphicx} %插入图片的宏包
% \usepackage{subfigure}
% \usepackage{float} %设置图片浮动位置的宏包

\usepackage{amsmath}
\usepackage{amssymb}
\usepackage{mathtools}

\usepackage{physics}

% \usepackage{siunitx}
% % 关闭警告
% \ExplSyntaxOn
% \msg_redirect_name:nnn { siunitx } { physics-pkg } { none }
% \ExplSyntaxOff

%%%%%%%%%%%%%%%%%%%%%%%% tcolorbox %%%%%%%%%%%%%%%%%%%%%%%%%%
\usepackage[most]{tcolorbox}

\newenvironment{tipbox}
{\begin{tcolorbox}[colback=green!5!white,colframe=green!75!black]}
{\end{tcolorbox}}

\newenvironment{conceptbox}
{\begin{tcolorbox}[colback=blue!5!white,colframe=blue!75!black]}
{\end{tcolorbox}}

\newenvironment{thmbox}
{\begin{tcolorbox}[colback=yellow!5!white,colframe=yellow!75!black]}
{\end{tcolorbox}}

\newenvironment{definition box}
{\begin{tcolorbox}[colback=blue!5!white,colframe=blue!75!black,parbox=false,before upper=\par,before lower=\par]}
{\end{tcolorbox}}

\newenvironment{theorem box}
{\begin{tcolorbox}[colback=yellow!5!white,colframe=yellow!75!black,parbox=false,before upper=\par,before lower=\par]}
{\end{tcolorbox}}
%%%%%%%%%%%%%%%%%%%%%%%%%%%%%%%%%%%%%%%%%%%%%%%%%%%%%%%%%%%%
            
%%%%%%%%%%%%%%%%%%%%% amsthm %%%%%%%%%%%%%%%%
\usepackage{amsthm}

\theoremstyle{definition}
\newtheorem{definition inner}{Definition}[section]
\newenvironment{definition}[1][]
{\begin{definition box}\begin{definition inner}[#1]\hfill\par}
{\end{definition inner}\end{definition box}}

\theoremstyle{plain}
\newtheorem{theorem inner}{Theorem}[section]
\newenvironment{theorem}[1][]
{\begin{theorem box}\begin{theorem inner}[#1]\hfill\par}
{\end{theorem inner}\end{theorem box}}

\theoremstyle{definition}
\newtheorem*{proof inner}{\textit{Proof}}
\renewenvironment{proof}[1][]
{\begin{proof inner}[#1]\hfill\par}
{\qed\end{proof inner}}

\theoremstyle{remark}
\newtheorem*{remark}{\textit{Remark}}
%%%%%%%%%%%%%%%%%%%%%%%%%%%%%%%%%%%%%%%%%%%%%%


\newcommand{\mycases}[2][0.8]{
    \left\{\vphantom{\scalebox{#1}{
        $\begin{matrix}#2\end{matrix}$
    }}\right.\!\!\!\!
    \begin{array}{l@{\quad}l}#2\end{array}
}

\newcommand{\e}{\mathrm{e}{\mkern 1 mu}}

\newcommand{\perm}[2]{\ensuremath{{}_{#1}\mathrm{P}_{#2}}}  % 排列数
\newcommand{\comb}[2]{\ensuremath{{}_{#1}\mathrm{C}_{#2}}}  % 组合数

\newcommand{\E}[1]{\mathbb{E}\qty[#1]}  % 期望
\newcommand{\Prob}{\mathbb{P}\qty}
\newcommand{\Var}[1]{\mathrm{Var}\qty(#1)}  % 方差

\begin{document}

\part*{Discrete Distributions}

\section{Random Variables of the Discrete Type}

\begin{definition}[Random Variables]
    Given a random experiment with sample space $S$,
    a function $X\colon S\rightarrow \mathbb{R}$
    that assigns a real number $X(s)=x$ to each $s\in S$
    is called a \textbf{random variable} (RV).
\end{definition}

\begin{definition}[Discrete Random Variable]
    A random variable is said to be of the \textbf{discrete type}
    if its range is \emph{finite} or \emph{countably infinite}.
\end{definition}

\begin{definition}[Probability Mass Function]
    A function $f\colon\mathbb{R}\rightarrow[0,1]$
    is called the \textbf{probability mass function} (pmf) of
    a DRV $X$ with range $S$ if it satisfies:
    \begin{itemize}
        \item $f(x)>0$ for $x\in S$, and $f(x)=0$ for $x\notin S$
        \item $\displaystyle\sum_{x\in S}f(x)=1$
        \item $\displaystyle \Prob(X\in A)=\sum_{x\in A}f(x)$, where $A\subseteq S$
    \end{itemize}
    Then $S$ is called the \textbf{support} or \textbf{space} of $X$.
\end{definition}

\begin{remark}
    From the third property, we have:
    \begin{equation*}
        \Prob(X=x) = f(x)
    \end{equation*}
\end{remark}


\begin{definition}[Cumulative Distribution Function]
    The function $F\colon\mathbb{R}\rightarrow[0,1]$
    defined by
    \begin{equation*}
        F(x) := \Prob(X\leq x)
    \end{equation*}
    is called the \textbf{cumulative distribution function} (cdf) of $X$.
\end{definition}
Then we have:
\begin{equation*}
    \Prob(a<X\leq b) = F(b) - F(a)
\end{equation*}

\begin{definition}[Line Graph and Probability Histogram]
    \centering
    \includegraphics*[width=0.5\textwidth]{assests/Line Graph and Histogram Graph.png}
\end{definition}


\begin{definition}[Uniform Distribution]
    An RV $X$ is said to have a \textbf{uniform distribution} if
    its pmf $f(x)$ is constant for all $x\in S$.
\end{definition}

\section{Mathmetical Expectation}

\begin{definition}[Mathmetical Expectation]
    Let $f(x)$ be the pmf of the DRV $X$ with space $S$.
    Suppose $g$ is a function of $X$, then
    \begin{equation*}
        \E{g(X)} := \sum_{x\in S}g(x)f(x)
    \end{equation*}
    is called the \textbf{mathmetical expectation} or \textbf{expected value} of $g$,
    if the sum exists.
    \begin{remark}
        $g(X)$ gives another RV.
    \end{remark}
\end{definition}

\subsection{Properties}

\begin{theorem}
    \begin{itemize}
        \item For any constant $c$:\quad $\E{c} = c$
        \item Linearity:\quad $\E{a X+b Y} = a\E{X} + b\E{Y}$
    \end{itemize}
\end{theorem}
\begin{proof}[Linearity]
    The linearity dose \textbf{not} depend on the independence of $X$ and $Y$.
    Here gives a proof for the continuous case.
    \begin{align*}
        \E{a X+Y} &=
    \int_{-\infty}^{\infty}\int_{-\infty}^{\infty} (a x+y) f_{X,Y}(x, y)\dd{x}\dd{y} \\
    &= a \int_{-\infty}^{\infty}x \int_{-\infty}^{\infty}f_{X,Y}(x, y)\dd{y}\dd{x}
    + \int_{-\infty}^{\infty}y \int_{-\infty}^{\infty}f_{X,Y}(x, y)\dd{x}\dd{y} \\
    &= a \int_{-\infty}^{\infty}x f_X(x)\dd{x}
    + \int_{-\infty}^{\infty}y f_Y(y)\dd{y} \\
        &= a \E{X} + \E{Y}
    \end{align*}
\end{proof}

\section{Special Mathmetical Expectations}

\subsection{Mean \& Variance}

\begin{definition}[Mean]
    The \textbf{mean} of a DRV $X$ (or of its distribution) is defined as
    \begin{equation*}
        \mu := \E{X} = \sum_{x\in S}x f(x)
    \end{equation*}
\end{definition}

\begin{definition}[Variance \& Standard Deviation]
    The \textbf{variance} of a DRV $X$ (or of its distribution) is defined as
    \begin{equation*}
        \sigma^2 := \Var{X} := \E{(X-\mu)^2} = \E{(X-\E{X})^2}
    \end{equation*}

    The \textbf{standard deviation} is defined as
    \begin{equation*}
        \sigma := \sqrt{\Var{X}}
    \end{equation*}
\end{definition}

\begin{remark}
    {Standard deviations} are not {mathmetical expectations}.
\end{remark}

\begin{remark}
    $b=\mu$ minimizes $\E{(X-b)^2}$.
    \begin{equation*}
        \Var{X} = \min_{b}\qty{ \E{(X-b)^2} }
    \end{equation*}
\end{remark}

\subsection{Properties of Variances}

\begin{theorem}
    Ways to compute the variance:
    \begin{equation*}
        \Var{X} = \E{(X-\mu_X)^2} = \E{X^2} - \E{X}^2 = \E{X(X-\mu_X)}
    \end{equation*}
\end{theorem}

\begin{theorem}
    \begin{itemize}
        \item $\Var{c} = 0$
        \item $\Var{c X} = c^2 \Var{X}$
    \end{itemize}
\end{theorem}

\subsection{The \texorpdfstring{$r$}{r}th Moment}

\begin{definition}[$r$th Moment]
    Given a positive integer $r$, the $r$th moment of an RV $X$ is defined as
    \begin{equation*}
        \E{X^r} = \sum_{x\in S}x^r f(x)
    \end{equation*}

    The $r$th moment of the distribution about $b$:
    \begin{equation*}
        \E{(X-b)^r} = \sum_{x\in S}(x-b)^r f(x)
    \end{equation*}

    The $r$th factorial moment:
    \begin{equation*}
        \E{(X)_r} :=  \E{X(X-1)(X-2)\cdots(X-r+1)}
    \end{equation*}
\end{definition}

\begin{definition}[Moment Generating Function]
    The \textbf{moment generating function} (mgf) of a DRV $X$ is defined as
    \begin{equation*}
        M_X(t) := \E{\e^{t X}} = \sum_{x\in S}\e^{t x} f(x)
    \end{equation*}
    if the sum exists and finite for $-h < t < h$.
\end{definition}

\begin{theorem}[Properties of Mgf]
    \begin{itemize}
        \item $M(0) = 1$
        \item $\displaystyle M^{(r)}(t) = \sum_{x\in S}x^r \e^{tx} f(x)$
        \item $\displaystyle M'(0) = \E{X}$
        \item $\displaystyle M^{(r)}(0) = \E{X^r}$
    \end{itemize}
\end{theorem}

\begin{remark}
    If moments of all order of $X$ exists, then
    \begin{equation*}
        M_X(t) = \sum_{k=0}^{\infty} \frac{\E{X^k}}{k!} t^k
    \end{equation*}
\end{remark}

\begin{theorem}
    2 RVs have the \textbf{same mgf}
    $\quad\iff\quad$
    they have the \textbf{same probability distribution},
    i.e., the \textbf{same pmf}
\end{theorem}


\section{The Binomial Distribution}


\subsection{The Bernoulli Distribution}

\paragraph{Bernoulli Experiment}
The outcomes can be classified in one of two mutually
exclusive and exhaustive ways, say either success or failure.

\begin{definition}[Bernoulli Distribution]
    Let $X$ be an RV associated with a \emph{Bernoulli experiment} with the
probability of success $p$.
    Then $X\in\{0,1\}$, and the pmf of $X$ is given by
    \begin{equation*}
        f(x) = p^x (1-p)^{1-x},\quad x\in\{0,1\}
    \end{equation*}
    Then we say $X$ has a \emph{Bernoulli distribution} with probability of success $p$.
\end{definition}


\begin{theorem}
    \begin{itemize}
        \item $\E{X} = p$
        \item $\Var{X} = p(1-p)$
        \item $M(t) = \e^t p + (1-p)$
    \end{itemize}
\end{theorem}


\begin{definition}[Bernoulli Trials]
    If a Bernoulli experiment is performed $n$ times
    \begin{itemize}
        \item independently
        \item with the probability of success remaining the same each
    \end{itemize}
    then these $n$ repetitions of the Bernoulli experiment is called
    $n$ Bernoulli trials.
\end{definition}

\paragraph{
    A random sample of size $n$ from a Bernoulli distribution
}
An observed sequence of $n$ Bernoulli trials that will be an $n$-tuple of zeros
and ones.


\subsection{The Binomial Distribution}

\paragraph{Desciption}
The DRV $X$ is the number of successes in $n$ Bernoulli trials.

\begin{definition}[Binomial Distribution]
    An RV $X$ is said to have a \textbf{binomial distribution},
    if the pmf is given by
    \begin{equation*}
        f(x) = \binom{n}{p}p^x (1-p)^{n-x} \qfor x\in\{0,1,2,\cdots,n\}
    \end{equation*}
    Then it is denoted by $X \sim \mathrm{B}(n,p)$.
\end{definition}

\begin{theorem}
    \begin{itemize}
        \item $M(t) = [\e^t p + (1-p)]^n$
        \item $\E{X} = n p$
        \item $\Var{X} = n p(1-p)$
    \end{itemize}
\end{theorem}


\section{The Negative Binomial Distribution}

\subsection{The Negative Binomial Distribution}

\paragraph{Description}
The DRV $X$ is the number of Bernoulli trials needed to get $r$ successes.

\begin{definition}[Negative Binomial Distribution]
    An RV $X$ is said to have a \textbf{negative binomial distribution},
    if the pmf is given by
    \begin{equation*}
        f(x) = \binom{x-1}{r-1}p^r (1-p)^{x-r} \qfor x\in\{r,r+1,\cdots\}
    \end{equation*}
\end{definition}


\begin{definition}[Negative Binomial Series]
    The negative binomial series is given by Maclaurin's series expansion:
    \begin{equation*}
        (1-w)^{-r}
        = \sum_{k=0}^{\infty} \binom{k+r-1}{r-1} w^k
        = \sum_{x=r}^{\infty} \binom{x-1}{r-1} w^{x-r}
        \qfor w\in(-1, 1)
    \end{equation*}
\end{definition}

\begin{theorem}
    \begin{itemize}
        \item $\E{X} = \dfrac{r}{p}$
        \item $\Var{X} = \E{X^2} - \E{X}^2 = \dfrac{1-p}{p^2}r$
        \item $\displaystyle
            M(t) = \E{\e^{t X}} = \left[\frac{p \e^t}{1-(1-p)\e^t}\right]^r
            \qfor (1-p)\e^t < 1$
    \end{itemize}
\end{theorem}

\subsection{The Geometric Distribution}

\paragraph{Description}
The DRV $X$ has a \emph{negative binomial distribution} with $r=1$.

\begin{definition}[Geometric Distribution]
    An RV $X$ is said to have a \textbf{geometric distribution} with the probability of success $p$,
    if the pmf is given by
    \begin{equation*}
        f(x) = p(1-p)^{x-1} \qfor x\in\{1,2,\cdots\}
    \end{equation*}
\end{definition}

\begin{theorem}
    For a positive integer $k$,
    \begin{gather*}
        \Prob(X > k) = \sum_{x=k+1}^{\infty} p(1-p)^{x-1} = (1-p)^k \\
        \Prob(X \leqslant k) = 1 - (1-p)^k
    \end{gather*}
\end{theorem}


\section{The Poisson Distribution}

\begin{definition}[Approximate Poisson Process (APP)]
    Let the number of occurrences of some event in a given continuous interval be counted.
    Then we have an APP with parameter $\lambda > 0$ if
    \begin{itemize}
        \item The number of occurrences in non-overlapping subintervals are \textbf{independent}.
        \item The probability of exactly one occurrence in a
        sufficiently short subinterval of length $h$ is
        approximately $\lambda h$.
        \item The probability of two or more occurrences in a
        sufficiently short subinterval is essentially $0$.
    \end{itemize}
\end{definition}

\paragraph{Description}
Let $X$ be the \textbf{number of occurrences} in an APP
within a \textbf{unit interval}.

\begin{remark}
    It can be seen as a binomial distribution
    with $n\to\infty$ and $p=\dfrac{\lambda}{n}$,
    i.e., $\displaystyle \mathrm{Poisson}(\lambda) = 
    \lim_{n\to\infty} \mathrm{B}(n,\frac{\lambda}{n})$.
\end{remark}

\begin{definition}[Poisson Distribution]
    An RV $X$ is said to have a \textbf{Poisson distribution} with parameter $\lambda$,
    if the pmf is given by
    \begin{equation*}
        f(x) = \frac{\lambda^x \e^{-\lambda}}{x!} \qfor x\in\{0,1,2,\cdots\}
    \end{equation*}
    Then it is denoted by $X \sim \mathrm{Poisson}(\lambda)$.
\end{definition}

\begin{theorem}
    \begin{itemize}
        \item $\displaystyle M(t) = \e^{\lambda (\e^t - 1)}$
        \item $\displaystyle \E{X} = \lambda$
        \item $\displaystyle \Var{X} = \lambda$
    \end{itemize}
\end{theorem}

The Poisson distribution works on a \textbf{unit interval},
so when considering an interval of length $T$,
we should also treat it as a unit interval.
Let $X$ be the number of occurrences.
Then $\E{X} = \lambda T$, $X\sim \mathrm{Poisson}(\lambda T)$.
\begin{equation*}
    f(x) = \frac{(\lambda T)^x \e^{-\lambda T}}{x!}
    \qfor x\in\{0,1,2,\cdots\}
\end{equation*}


\end{document}
